\chapter{Introduction}

\label{ch:intro}

  

  

\section{Background}

  

  

  

\subsection{OAM and Future Optical Communications}

  

傳統上，光通訊主要依賴空間、時間、波長、振幅與偏振等自由度來調變訊號。然而，隨著全球數據流量迎來爆炸性成長，傳統自由度所組成的多工系統在容量提升上，逐漸面臨頻譜效率、功耗與系統複雜度等實際工程限制。因此，探索新型態的多維度光場調變技術，成為下一代光通訊技術發展的重要方向。

  

在這些新興的光場自由度中，軌道角動量 (Orbital Angular Momentum, OAM) 以其獨特的螺旋相位結構受到廣泛關注。攜帶 OAM 的光束相位會隨方位角改變(方位角相位依賴性，以拓樸電荷 $\ell$ 表示，在理論上可取任意整數值。由於不同拓樸電荷的 OAM 模態之間彼此正交，能夠提供空間複用額外的空間模態自由度，在不額外佔用頻譜資源的前提下提升頻譜效率。

  

這種在自由空間光通訊與無線傳輸的 OAM 技術，已在諸多文獻中獲得廣泛驗證；其中 Willner 等人對 OAM 通訊技術進行的文獻回顧，更奠定了 OAM 在下一代高容量光學通訊中的關鍵地位\cite{willner2015optical}。除了經典光通訊外，OAM 光束在高維度量子資訊處理等前沿光子學領域中，亦展現出極高的應用價值。

  

  

\subsection{Need for Chip-scale OAM Emitters}

  

儘管 OAM 光束的生成技術已相當多元，包括螺旋相位片、空間光調製器及電腦生成全息圖等，但此類方法多依賴體積龐大且難以微縮的自由空間光學元件。這不僅使主動式元件的響應速度受限，更難以整合緊湊型光子積體電路及晶片上光源。

  

超穎表面的出現為晶片級光場調控提供了突破性的解決方案。透過次波長奈米結構陣列對入射光的振幅、相位及偏振進行精確的局部調控，超穎表面得以在數百奈米厚度的二維介面上實現傳統巨觀光學元件的功能。其中，幾何相位超穎表面因其簡單的相位編碼機制與良好的製程容忍度，成為生成 OAM 光束最廣泛採用的晶片級路徑之一。

  

然而，超穎表面本質上屬於被動波前轉換器，必須依賴外部光源入射，無法主動產生具備手性的光場。因此，如何引入具備天然發光能力的主動材料，並將其與奈米結構進行晶片級的混合式整合，成為實現晶片上主動式手性光源的關鍵挑戰。

  

\subsection{Active Chiral Light Sources}

  

若要實現主動式的晶片光源，必須引入具備天然手性發光能力的主動材料。單層二維過渡金屬硫屬化物 (Transition Metal Dichalcogenides, TMDs，如 MoS$_2$、WS$_2$ 等) 是一類具備直接能隙的二維半導體。由於其獨特的晶體結構，這類材料展現出特殊的谷光學特性，在內部電子狀態與發光偏振態之間建立了直接聯繫。這種特性使單層 TMDs 能在特定光激發或電注入下，產生具備特定旋向的圓偏振光，為新型光電元件的開發提供了天然平台，因而成為近年發展固態手性光源的研究熱點。

  

\subsection{Toward Active Valley-controlled OAM Sources}

  

將單層 TMDs 的谷極化主動發光與緊湊型整合的超穎表面相結合，為晶片級自旋角動量 (Spin Angular Momentum, SAM) 至 OAM的轉換提供了完整的物理連接。理想情況下，單層 TMDs 的谷極化發光攜帶特定的 SAM，隨後透過緊湊型整合的超穎表面進行幾何相位調控，即可將其轉化為特定的 OAM 模態。

  

這種主動材料與被動奈米結構的混合式整合架構，不僅繼承了二維材料獨特的谷電子學自由度，亦發揮了超穎表面在亞波長尺度下極致的波前操縱能力，是實現高整合度、主動可調控晶片級 OAM 發射器的理想演進方向。

  

\section{Motivation}

儘管前述各小節已揭示 OAM 光束的應用潛力與單層 TMDs 的具備的自旋–谷鎖定特性與手性發光機制，在晶片級實現上，至今仍面臨發光效率、訊號串擾、元件整合以及波前控制能力等多重挑戰。

\subsection{Limitation of Existing TMD-OAM Platforms}

就元件整合而言，傳統上產生或調變 OAM 光束多依賴空間光調變器 (SLMs) 或叉形全息光柵等毫米級元件。這些巨觀元件體積龐大且需要複雜的實驗光路對準，難以滿足下一代光電晶片對於積體化的迫切需求。若將 TMDs 與超薄金屬電漿子超穎表面結合以縮減體積，雖然能將厚度壓縮至次波長尺度，但金屬的高歐姆損耗嚴重限制能量轉換效率，使得僅靠元件結構的縮減仍難以突破效能瓶頸。究其原因，除了結構損耗，更深層的內在物理限制在於單層 TMDs 的激子發光極易受到內部谷間散射與去極化機制的影響，導致有限的谷極化純度。這種本質上的去極化發光特性，使得後續在晶片尺度上進行 OAM 編碼的保真度受到嚴重限制，難以實現高對比度的通道分離。

  

\subsection{Limitations of Conventional Metasurface Design}

為了克服上述電漿子元件的損耗問題，利用全介電質超穎表面進行波前調控已成為主流，然而傳統的超穎表面設計方法學面臨本質上的物理限制。就設計方法論而言，傳統正向設計依賴物理直覺建構幾何單元，其設計自由度局限於少數幾何參數。面對需同時調控振幅、相位、偏振與 OAM 模態純度的多目標要求時，正向設計難以在龐大參數空間中尋得全域最佳解。此種設計方式亦面臨空間解析度與近場耦合之間的權衡。當正向設計縮小結構間距以獲取精細相位梯度時，更會因相鄰結構的近場耦合而引入相位雜訊並限制繞射效率。因此，解決上述問題的關鍵在於採用能釋放全空間幾何自由度的反向設計架構。

\section{Research Objectives and Contributions}

  

本研究的核心目標在於建立一個基於反向設計的全介電質超穎表面平台，探索單層過渡金屬硫屬化物谷極化發光至高純度軌道角動量光束之轉換機制。

  

為達成此目標，本工作首先建構整合單層過渡金屬硫屬化物發光體、全介電質超穎表面以及近遠場分析之三維數值模擬架構。在此平台基礎上，進一步發展基於伴隨變量法的拓撲優化方法，以有效拓展空間幾何自由度，進行谷控制軌道角動量發射器之反向設計。

  

於優化過程中，將遠場拉格耳–高斯模態純度與數值模擬與理想目標的強度重疊積分納入核心效能指標，以評估並提升所生成 OAM 模態之品質。最後，透過最佳化結構分析其自旋–軌道轉換與模態形成機制，建立由數值設計、性能評估至物理機制分析之完整研究流程。

  

\section{Thesis Organization}

  

本論文架構如下：

  

第一章引介研究背景、研究動機、研究目標、研究貢獻以及全文的組織架構。

  

第二章系統性地回顧軌道角動量理論、單層過渡金屬硫屬化物的谷選擇性發光、基於超穎表面的軌道角動量產生技術，以及相關的反向設計方法。

  

第三章詳細闡述本研究所採用的數值模擬架構、物理模型、拓樸最佳化方法以及效能指標的具體定義。

  

第四章展示最佳化元件結構，並透過近場與遠場特徵分析評估其光學效能。

  

第五章探討軌道角動量模態生成的物理機制，將所得結果與前人研究進行對比，並剖析該方法的物理限制與未來研究方向。

  

第六章總結本論文的主要研究結論，並概述未來潛在的研究方向。
